%%
%% fall-lecture19.tex
%% 
%% Made by Alex Nelson <pqnelson@gmail.com>
%% Login   <alex@lisp>
%% 
%% Started on  2026-11-10T09:51:08-0800
%% Last update 2026-11-10T09:51:08-0800
%% 

\lecture{}

We will continue determining relevant physical quantities for quasi 1d flows.

\begin{notation}
The notation for the speed of sound is sometimes $c$, other times it
is $a$. This is unfortunate, but such is life.
\end{notation}

\begin{node}[Speed of sound]
Now, recall the definition for the speed of sound is
\begin{equation}
a^{2} = \left(\frac{\partial p}{\partial\rho}\right)_{S}.
\end{equation}
If the flow is isentropic [which happens if $f=0$ and $\dot{q}=0$],
then we have the difference equation 
\begin{equation}
\D p = a^{2}\,\D\rho.
\end{equation}
Then we can use the conservation of momentum $u\partial_{x}u+(1/\rho)\partial_{x}p=0$
[hence $-\rho u\,\D u =\D p$] to get
\begin{equation}
-\rho u\,\D u = a^{2}\,\D\rho.
\end{equation}
This gives us, using the Mach number $\Mach^{2}=u^{2}/a^{2}$,
\begin{equation}
\frac{\D\rho}{\rho}=\frac{-u}{a^{2}}\D u =\frac{-u^{2}}{a^{2}}\frac{\D u}{u}=-\Mach^{2}\frac{\D u}{u}.
\end{equation}
Then the relative change in density is given by
\begin{equation}\label{eq:fall:lec19:me101a:eq1}
\boxed{\frac{\D\rho}{\rho}=-\Mach^{2}\frac{\D u}{u}.}
\end{equation}
\end{node}

\begin{node}
The other equation, by continuity [conservation of mass],
\begin{equation}
\D(\rho u A)=0,
\end{equation}
which expands to (via the product/Leibniz rule),
\begin{equation}
\rho u\,\D A + \rho A\,\D u + Au\,\D\rho=0.
\end{equation}
Then dividing through by $\rho u A$ gives us
\begin{equation}\label{eq:fall:lec19:me101a:eq2}
\boxed{\frac{\D A}{A} + \frac{\D u}{u} + \frac{\D\rho}{\rho}=0.}
\end{equation}
\end{node}

\begin{node}
Then combining Eq~\eqref{eq:fall:lec19:me101a:eq1}
and Eq~\eqref{eq:fall:lec19:me101a:eq2}, we obtain:
\begin{equation}\label{eq:fall:lec19:me101a:eq3}
  \boxed{\frac{\D u}{u}=\frac{1}{\Mach^{2}-1}\frac{\D A}{A}.}
\end{equation}
\end{node}

\begin{proof}
We see that Eq~\eqref{eq:fall:lec19:me101a:eq2} gives us,
\begin{equation}
\frac{\D\rho}{\rho} = -\frac{\D u}{u} -\frac{\D A}{A}.
\end{equation}
Then noting this is the left-hand side of Eq~\eqref{eq:fall:lec19:me101a:eq1},
we find
\begin{subequations}
  \begin{align}
-\Mach^{2}\frac{\D u}{u} &= \frac{\D\rho}{\rho}\quad\mbox{by Eq~\eqref{eq:fall:lec19:me101a:eq1}}\\
&= -\frac{\D u}{u} -\frac{\D A}{A},
  \end{align}
\end{subequations}
then adding $(\D u)/u$ to both sides yields
\begin{equation}
(1-\Mach^{2})\frac{\D u}{u} = -\frac{\D A}{A}.
\end{equation}
Dividing through by $1-\Mach^{2}$ gives us the result.
\end{proof}


\subsection{Conditions leading to flow acceleration}

\begin{node}
We want to find what conditions leads us to $u>0$ and $\D u>0$ (or, if
the flow is ``moving to the left'', $u<0$ and $\D u<0$). We will
divide up our attention into various ``regimes'' of the Mach number.
\end{node}

\begin{case}[$\Mach=0$]\label{case:fall:lec19:me101a:case1}
Then by Eq~\eqref{eq:fall:lec19:me101a:eq1} we have
\begin{equation}
\D\rho=0.
\end{equation}
This is the incompressible flow regime. And Eq~\eqref{eq:fall:lec19:me101a:eq3}
gives us
\begin{equation}
\frac{\D u}{u}=-\frac{\D A}{A}>0,
\end{equation}
which implies (since $(\D u)/u>0$ and $A>0$):
\begin{equation}
\D A < 0.
\end{equation}
\end{case}

\begin{case}[$0<\Mach<1$---the Subsonic regime]
When $0<\Mach<1$, then $1-\Mach^{2}>0$. Then Eq~\eqref{eq:fall:lec19:me101a:eq3}
gives us
\begin{equation}
\D u>0\implies\D\rho<0.
\end{equation}
So the density is not constant, and its change $\D\rho$ depends on the
Mach number.
\end{case}

\begin{case}[$1<\Mach$---the Supersonic regime]
When we are in the ``supersonic regime'' $1<\Mach$, then $1-\Mach^{2}<0$
and Eq~\eqref{eq:fall:lec19:me101a:eq3} gives us
\begin{equation}
\D u>0\implies\D A>0.
\end{equation}
Eq~\eqref{eq:fall:lec19:me101a:eq1} tells us, since $\D u>0$, we must
have $\D\rho<0$---we are in a compressible flow.
\end{case}

\begin{case}[$\Mach=1$---the Sonic regime]
When $\Mach=1$, Eq~\eqref{eq:fall:lec19:me101a:eq3} becomes ill-defined
\begin{equation}
\frac{\D u}{u}=\frac{1}{0}\frac{\D A}{A},
\end{equation}
which means if $(\D u)/u$ is finite, then we must have
\begin{equation}
\frac{\D A}{A}=0.
\end{equation}
This means $\D A=0$. So we have two possibilities:
\begin{center}% from https://tex.stackexchange.com/a/17102/14751
(1)\begingroup\setbox0=\hbox{\includegraphics{img/img.40}}\parbox{\wd0}{\box0}\endgroup
\quad and\quad(2)
\begingroup\setbox0=\hbox{\includegraphics{img/img.41}}\parbox{\wd0}{\box0}\endgroup
\end{center}
But (2) does not make sense in light of Case~\eqref{case:fall:lec19:me101a:case1},
because the fluid leaves this muzzle in the subsonic regime.
\end{case}

\begin{note}
For a non-classic gas dynamics regime, $\Gamma<0$, we could
potentially have a flow passing through a nozzle which looks like the
situation we eliminated. But this is a highly theoretical concern
which we will not encounter in practice.
\end{note}

\begin{note}
For the physically realistic sonic regime, $\Gamma>1$ for ideal and
$0<\Gamma<1$ for non-ideal flows (where  $\Gamma$ is the fundamental
derivative for gas dynamics).
\end{note}

\subsection{Material property $\beta$}

\begin{definition}
We define a new quantity $\beta$ in one of two equivalent ways:
\begin{enumerate}
\item $\beta :=\Gamma-1$ (or $\Gamma = 1+\beta$), or
\item $\displaystyle{\beta:=\frac{\rho}{a}\left(\frac{\partial a}{\partial\rho}\right)_{S}=\rho a\left(\frac{\partial a}{\partial p}\right)_{S}}$
\end{enumerate}
Thompson~\cite[p.1849~\S VII]{thompson1971fundamental} calls $2(\Gamma-1)$ the
\define{Parameter of Nonlinearity} (and I have yet to see anyone else
give $\beta$ a name).
\end{definition}

\begin{node}
These two definitions for $\beta$ are equivalent, since [proving (2)
  gives us (1)]:
\begin{subequations}
  \begin{align}
\Gamma &= \frac{V^{3}}{2a^{2}}\left(\frac{\partial^{2}p}{\partial V^{2}}\right)_{S}\\
&= \frac{a^{4}}{2V^{3}}\left(\frac{\partial^{2}V}{\partial p^{2}}\right)_{S}\\
&= \dots\\
&= 1 + \frac{\rho}{a}\left(\frac{\partial a}{\partial \rho}\right)_{S}=1+\beta.
  \end{align}
\end{subequations}
We can read these equalities from bottom to top to prove (1) gives us (2).
\end{node}

\begin{node}
If we're dealing with an isentropic process $\D S=0$,
then
\begin{subequations}
  \begin{align}
\frac{\D a}{a} &= \frac{\rho}{a}\frac{\D\rho}{\rho}\frac{\D a}{\D\rho}\\
&= \frac{\rho}{a}\frac{\D a}{\D\rho}\frac{\D\rho}{\rho}\\
&=\beta\frac{\D\rho}{\rho}\quad\mbox{along isentrope}.
  \end{align}
\end{subequations}
This tells us the relative change in the speed of sound $a^{-1}\,\D a$
is directly proportional to the relative change in density $\rho^{-1}\,\D\rho$
with proportionality constant given by $\beta$.


We can write this further as
\begin{equation}
\frac{\D a}{a} = \beta\frac{\D p}{\rho a^{2}}.
\end{equation}
So, 
\begin{equation}
\frac{\D\rho}{\rho}=-\Mach^{2}\frac{\D u}{u}\implies\frac{\D a}{a}=-\beta\Mach^{2}\frac{\D u}{u}.
\end{equation}
Since $\Mach=u/a$, we find
\begin{subequations}
  \begin{align}
\frac{\D\Mach}{\Mach} &= \frac{1}{M}\left(\frac{\D u}{a}-\frac{u}{a^{2}}\D a\right)\\
&= \frac{a}{u}\frac{\D u}{a}-\frac{a}{u}\frac{u}{a^{2}}\D a\\
&=\frac{\D u}{u}-\frac{\D a}{a},
  \end{align}
\end{subequations}
hence we obtain (combining these steps all together):
\begin{equation}\label{eq:fall2026:me101a:lec19:change-in-u-related-to-mach-number}
\boxed{\frac{\D\Mach}{\Mach} = (1+\beta\Mach^{2})\frac{\D u}{u}.}
\end{equation}
We also obtain
\begin{equation}
\boxed{\frac{\D A}{A}=\frac{\Mach^{2}-1}{1+\beta\Mach^{2}}\frac{\D\Mach}{\Mach}.}
\end{equation}
\end{node}

\begin{node}
From previous relationships, we also obtain,
\begin{equation}
\boxed{\frac{\D\rho}{\rho}=\frac{-\Mach^{2}}{1+\beta\Mach^{2}}\frac{\D\Mach}{\Mach}.}
\end{equation}
\end{node}

\begin{node}
We are missing a relation expressing the change in energy in terms of
the change in the Mach number (we have boxed every other such relationship).
We see for enthalpy,
\begin{subequations}
  \begin{align}
\D h &=-u\,\D u\\
&=-u^{2}\frac{\D u}{u}\\
&=2(h-h_{0})\frac{\D u}{u},\\
\intertext{(where $h_{0}$ is total enthalpy) hence:}
\frac{\D u}{u} &= \frac{\D h}{2(h-h_{0})}.
  \end{align}
\end{subequations}
Then we can plug in the relation Eq~\eqref{eq:fall2026:me101a:lec19:change-in-u-related-to-mach-number}
to find
\begin{equation}
\boxed{\frac{\D h}{2(h-h_{0})}=\frac{1}{1+\beta\Mach^{2}}\frac{\D\Mach}{\Mach}.}
\end{equation}
We could have used the equation of state
\begin{equation}
\D h = T\,\D S + V\,\D p = 0 + V\,\D p
\end{equation}
where $\D S=0$ (since we're on an isentrope), which gives us
\begin{equation}
\left(\frac{\partial h}{\partial p}\right)_{S}=\frac{1}{\rho},
\end{equation}
which gives us
\begin{equation}\label{eq:fall2026:me101a:enthalpy-differential}
\boxed{\D h=\frac{1}{\beta}a\,\D a.}
\end{equation}
\end{node}

\subsection{Stagnation Quantities}

\begin{node}
If $\beta$ is constant, then we can integrate these boxed
relationships. For example, Eq~\eqref{eq:fall2026:me101a:enthalpy-differential}
gives us
\begin{equation}
\int^{h}_{h_{0}}\D h'=\int^{a}_{a_{0}}\frac{a'}{\beta}\,\D a'
\end{equation}
where $a_{0}$ is the reference speed of sound and $h_{0}$ is the total
enthalpy (not a typo). Taking $a_{0}=0$, this integral gives us
\begin{equation}
\boxed{h = \frac{1}{2\beta}a^{2}.}
\end{equation}
Similarly, the other relationships
\begin{equation}
\int^{a}_{a_{0}}\frac{\D a'}{a'}=\int^{\rho}_{\rho_{0}}\beta\frac{\D\rho'}{\rho'}\implies
\frac{a}{a_{0}}=\left(\frac{\rho}{\rho_{0}}\right)^{\beta},
\end{equation}
\begin{equation}
\int^{\rho}_{\rho_{0}}\frac{\D\rho'}{\rho'}=-\int^{\Mach}_{0}\frac{(\Mach')^{2}}{1+\beta\Mach'}\frac{\D\Mach'}{\Mach'}\implies\frac{\rho}{\rho_{0}}=(1+\beta\Mach^{2})^{1//(2\beta)}.
\end{equation}
The values $\rho_{0}$, $a_{0}$, $p_{0}$, \dots, (all subscripted with
zero) are called the \define{Stagnation Quantities}. They describe the
situation where all kinetic energy has been transformed into internal
energy. For any flow, we can move along an isentropic process to a
state where $u=0$, $\Mach=0$, $\rho=\rho_{0}$, $a=a_{0}$, and it is
this stagnation point.
\end{node}

\begin{aside}[Stagnation ondition]
The stagnation quantities (with subscript 0) $T_{0}$, $u_{0}$,
$p_{0}$, $\rho_{0}$, etc., are all achieved with a flow decelerated
isentropically [i.e., no heat loss through $Q$ or friction, and no
  shock waves].
We move along isentropes. For compressible fluids, this becomes very
important and we define stagnation quantities as our ``base line''. We
achieve this in two ways:
\begin{enumerate}
\item Using \define{Pitot tubes} which are designed to slow supersonic
  fluids down isentropically to measure $p_{0}$, $\rho_{0}$, $T_{0}$
\item \define{Stagnation Chambers} where we start at the stagnation
  state (so we can measure $p_{0}$, $\rho_{0}$, $T_{0}$, etc.) then
  speed up the fluid.
\end{enumerate}
\end{aside}

\begin{proposition}
We have $\displaystyle{\frac{p_{0}}{p}=(1+\beta\Mach^{2})^{(2\beta+1)/(2\beta)}}$.
\end{proposition}

\begin{proof}
Since $\D p = a^{2}\,\D\rho$, we find
\begin{subequations}
  \begin{align}
\int^{p}_{p_{0}}\frac{\D p'}{\rho_{0}a_{0}^{2}}
&=\int^{\rho}_{\rho_{0}}\frac{\rho'a^{2}}{\rho_{0}a_{0}^{2}}\frac{\D\rho'}{\rho'}\\
&=\int^{\Mach}_{0}\left(\frac{\rho}{\rho_{0}}\right)
\left(\frac{a}{a_{0}}\right)^{2}
\frac{-(\Mach')^{2}}{1+\beta(\Mach')^{2}}
\frac{\D\Mach'}{\Mach'},
  \end{align}
\end{subequations}
which gives us the relationship
\begin{equation}
\frac{p-p_{0}}{\rho_{0}a_{0}^{2}} 
= \frac{1}{2\beta+1}\left[(1+\beta\Mach^{2})^{-(2\beta+1)/(2\beta)}-1\right],
\end{equation}
which tells us
\begin{equation}
\lim_{\Mach\to0}p=0.
\end{equation}
This gives us
\begin{equation}
p_{0}=\frac{1}{2\beta+1}\rho_{0}a_{0}^{2},
\end{equation}
hence
\begin{equation}
\frac{p_{0}}{p}=(1+\beta\Mach^{2})^{(2\beta+1)/(2\beta)},
\end{equation}
as desired.
\end{proof}

\begin{proposition}
We have $\displaystyle{\frac{h_{0}}{h}=1+\beta\Mach^{2}}$.
\end{proposition}

\begin{node}
The upshot to all of this is: we want to relate quantities (like $p$,
$\rho$, $T$, etc.) to stagnation quantities (resp., $p_{0}$,
$\rho_{0}$, $T_{0}$, etc.) and the Mach number $\Mach$. These
relationships allow us to know only two things:
\begin{enumerate}
\item the values of the stagnation quantities, and
\item the current value of the Mach number.
\end{enumerate}
From these, we can immediately calculate the current values for the
physical quantities $p$, $\rho$, $T$, etc.
\end{node}

\begin{note}
We did not compute the cross sectional area for stagnation. That's
because there is no meaningful ``stagnation area $A_{0}$''! But there
is an $A^{*}$ where $\Mach=1$ is sonic (we call $A^{*}$ the
\define{Sonic Area}). Then:
\begin{subequations}
  \begin{equation}
\int^{A}_{A^{*}}\frac{\D A'}{A'}=\int^{\Mach}_{1}\frac{(\Mach')^{2}-1}{1+\beta(\Mach')^{2}}\frac{\D\Mach'}{\Mach'},
  \end{equation}
  hence
\begin{equation}
\left(\frac{A}{A^{*}}\right)^{2}=\frac{1}{\Mach^{2}}\left(\frac{1+\beta\Mach}{1+\beta}\right)^{(1+\beta)/\beta}.
\end{equation}
\end{subequations}
\end{note}

\begin{remark}[References]
For discussion about stagnation points, see Batchelor~\cite[\S2.7~p.105]{batchelor1973introduction},
Anderson~\cite[\S7.5]{anderson1984fundamentals}.
Oddly, Landau and Lifshitz~\cite[\S10~p.20~\textit{et~seq.}]{landau1987fluid} seem to just assume the reader is
familiar with stagnation points, and does not take advantage of
stagnation quantities.\footnote{Perhaps in \S114 (around p.458) or somewhere, Landau
and Lifshitz takes advantage of stagnation quantities?} (I presume it is because aerodynamics takes
advantage of stagnation points, with dividends.)
\end{remark}

\begin{remark}[TODO]
I should work out the physics of a Pitot tube, and a Stagnation Chamber.
\end{remark}